\subsection{Conservative Admission and the Queue-Free MDP Invariant}
\label{app:conservative_admission}
This optional model restricts charging candidates before assignment;
it is an alternative sufficient construction for
Assumption~\ref{ass:reservation}, not an exact description of all
physically feasible charging decisions.

\paragraph{State and reservations.}
Let $\mathcal C$ contain the stations subject to this admission rule,
with physical plugs $\mathcal P_c=\{1,\ldots,C_c\}$, $C_c\ge1$.
At decision time $t$, write the augmented MDP state as
$S_t=(Z_t,\mathcal R_t)$, where $Z_t$ contains the remaining vehicle,
request, and decision-phase information. The reservation state
$\mathcal R_t$ records each outstanding job $r$ by its vehicle,
station $c_r$, plug $p_r$, arrival $a_r$, and service interval
$I_r=[b_r,e_r)$, $e_r>b_r$.
For each plug, define the disjoint job sets
$\mathcal R^{\rm ch}_{c,p,t}$ (currently charging) and
$\mathcal R^{\rm in}_{c,p,t}$ (committed but not yet arrived), and
\begin{equation}
\begin{aligned}
\mathcal R_{c,p,t}&=\mathcal R^{\rm ch}_{c,p,t}
                     \mathbin{\dot\cup}\mathcal R^{\rm in}_{c,p,t},\\
H_{c,p,t}&=\bigcup_{r\in\mathcal R_{c,p,t}} I_r,\\
B_{c,t}(\tau)&=\sum_{p\in\mathcal P_c}
               \sum_{r\in\mathcal R_{c,p,t}}\mathbf1\{\tau\in I_r\}.
\end{aligned}
\label{eq:conservative_background}
\end{equation}
The calendar is complete and consistent with actual charging: every
vehicle charging or committed to arrive at a controlled station appears
exactly once. Let $\mathcal S_0$ be the states with
empty charging queues, pairwise disjoint reservations on every plug,
and $b_r=a_r$ for every committed inbound job. All arrival times and
positive service durations are determined by the state and candidate
action. Only new decision vehicles generate potential jobs; a committed
job is not also counted as a new candidate.

\paragraph{Virtual admission calendar.}
At station $c$, order all physical candidates
$k_1,\ldots,k_{m_c}$ by $(a_{k,c,t},k)$, where
$a_{k,c,t}=t+\Delta_{k,c}$ and $h_{k,c,t}>0$ is the full service duration.
Its immediate-service interval is
$I_{k,c,t}=[a_{k,c,t},a_{k,c,t}+h_{k,c,t})$.
Let $f_{k,c,t}\in\{0,1\}$ be the otherwise-applicable edge mask.
For this station suppress $c,t$ and write $a_i=a_{k_i,c,t}$,
$h_i=h_{k_i,c,t}$. Set $L_p^0=t$. Recursively define
\begin{equation}
\begin{aligned}
g_{i,p}&=\min\{s\ge\max(a_i,L_p^{i-1}):\\[-2pt]
&\qquad[s,s+h_i)\cap H_{c,p,t}=\varnothing\},\\
p_i&=\operatorname*{arg\,min}_{p\in\mathcal P_c}(g_{i,p},p),
\qquad \sigma_i=g_{i,p_i},\\
L_p^i&=\begin{cases}\sigma_i+h_i,&p=p_i,\\
                         L_p^{i-1},&p\ne p_i.
          \end{cases}
\end{aligned}
\label{eq:conservative_virtual_recursion}
\end{equation}
The pair in the argmin is ordered lexicographically. All candidates,
including rejected ones, retain their virtual intervals
$\widehat I_i=[\sigma_i,\sigma_i+h_i)$ while this calendar is built.
The conservative mask and charging selectors satisfy
\begin{equation}
\begin{aligned}
\chi_{k_i,c,t}&=f_{k_i,c,t}\mathbf1\{\sigma_i=a_i\},\\
y_{k,c,t}&\in\{0,1\},\qquad y_{k,c,t}\le\chi_{k,c,t},\\
\sum_{c\in\mathcal C}y_{k,c,t}&\le1\qquad\forall k.
\end{aligned}
\label{eq:conservative_mask}
\end{equation}
Write $\mathcal U_{\rm cons}(s)$ for joint actions obeying this mask
and all other action constraints. Only selected intervals become real
commitments; unselected virtual jobs are discarded after the solve.

\begin{proposition}[Conservative charging safety]
\label{prop:conservative_queue_safety}
For $S_t=s\in\mathcal S_0$, every
$u\in\mathcal U_{\rm cons}(s)$ admits immediate service for every
selected charging arrival, without delaying any existing reservation.
Suppose further that all arrivals to $\mathcal C$ obey this admission
rule, selected absolute service intervals remain commitments in later
decisions, and the transition executes those deterministic times,
releasing completed service before simultaneous arrivals. Then
\begin{equation}
\begin{aligned}
P(\mathcal S_0\mid s,u)&=1\qquad\forall s\in\mathcal S_0,\\
&\hspace{9mm}\forall u\in\mathcal U_{\rm cons}(s),
\end{aligned}
\label{eq:conservative_safe_kernel}
\end{equation}
where $P$ is the augmented MDP transition kernel. Consequently, for
any policy $\pi$ that is well defined for all decision epochs and
selects an admissible action at every reachable state, with
$S_0\in\mathcal S_0$,
\begin{equation}
\Pr^\pi\{Q_c(\tau)=0\ \forall c\in\mathcal C,
                              \ \forall\tau\ge0\}=1,
\label{eq:conservative_queue_invariant}
\end{equation}
where $Q_c(\tau)$ is the number of vehicles waiting for a plug.
This claim does not assert that an admissible joint action always exists.
\end{proposition}
\begin{proof}
For each plug, the recursion places a new virtual interval after all
earlier virtual intervals on that plug and outside every fixed
reservation. Induction over $i$ therefore gives
\begin{equation}
\sum_{r\in\mathcal R_{c,p,t}}\mathbf1\{\tau\in I_r\}
+\sum_{i:p_i=p}\mathbf1\{\tau\in\widehat I_i\}\le1
\quad\forall p,\tau.
\label{eq:conservative_plug_bound}
\end{equation}
A selected candidate has $\sigma_i=a_i$, so its actual service
interval is exactly its virtual interval. Deleting all unselected
virtual intervals and summing over plugs yields
\begin{equation}
B_{c,t}(\tau)+\sum_k y_{k,c,t}
\mathbf1\{\tau\in I_{k,c,t}\}\le C_c
\quad\forall c,\tau.
\label{eq:conservative_occupancy_bound}
\end{equation}
The retained per-plug allocation certifies immediate service and leaves
every previous commitment intact. Equal-time departures release their
plugs before arrivals, so this allocation is executable without a queue.
Recording selected intervals, discarding completed jobs, and preserving
all remaining commitments gives another state in $\mathcal S_0$.
All stochastic request or acceptance outcomes preserve this argument
provided their charging actions obey the same admission rule. Hence
\eqref{eq:conservative_safe_kernel} holds. Induction over decisions,
together with the full-interval bound between decisions, proves
\eqref{eq:conservative_queue_invariant}.
\end{proof}

\paragraph{Scope.}
Independent station calendars can reserve virtual capacity for the
same vehicle at several stations, although only one is selected.
Thus the rule is sufficient and potentially restrictive. Additional
group quotas only restrict this safe set further; flow and SSG
exactness refer to the resulting fixed assignment model. The MDP
guarantee requires the stated transition semantics, not just a mask
computed at one epoch. In particular, an arrival-before-release
implementation or an unrecorded autonomous arrival is not covered.
